\section{From conformal fields to explicit lattice states}
\label{sec:construction}
\subsection{What is being constructed and measured?}
A periodic chain contains $N$ two-level sites. A basis configuration is
$\ket{n_0\cdots n_{N-1}}$, with $n_j=0,1$. A wavefunction is a table of
complex amplitudes, $\ket{a}=\sum_{\bm n}\psi_a(\bm n)\ket{\bm n}$,
normalized so that $\sum_{\bm n}|\psi_a(\bm n)|^2=1$.
An interval $A$ consists of the first $\ell$ consecutive sites; its complement
is $\bar A$. We trace out $\bar A$ to form
\begin{equation}
 (\rho_A^a)_{uv}=\sum_w\psi_a(u,w)\psi_a(v,w)^*,\qquad x=\ell/N.
 \label{eq:rdm}
\end{equation}
Here $u,v$ label interval configurations, $w$ labels the complement, and
the star denotes complex conjugation. Throughout this report,
\begin{equation}
 m_{ab}=\tfrac12(\rho_A^a+\rho_A^b),\qquad
 \Dab=\Tr\rho_A^a(\log\rho_A^a-\log m_{ab}).
 \label{eq:observable}
\end{equation}
Logarithms are natural, so relative entropy is measured in nats. The
mixture probability is $p=1/2$ and $0\leq\Dab\leq\log2$. This is an
incoherent mixture of density matrices, not a coherent sum of state vectors.
The ordering $a\mid b$ identifies the numerator; reversing it gives a
different observable unless a specific symmetry proves equality.

Conformal field theory (CFT) describes the long-distance limit of a critical
chain. A \emph{primary field} creates a basic excitation; its left and right
scaling weights are $h,\bar h$, its scaling dimension is $h+\bar h$, and its
conformal spin is $h-\bar h$. Descendants are obtained by acting with modes
of the symmetry generators. A \emph{conformal block} is a correlation
amplitude with specified intermediate fusion channels, that is, specified
outcomes when fields are brought together. We turn these amplitudes into
wavefunction coefficients by associating the field components or fusion
channels with the physical basis configurations. This construction chooses
a lattice representative; it does not guarantee that every arbitrary field
insertion is an eigenstate of a chosen lattice Hamiltonian.

The production constructions are given in Secs.~\ref{sec:bosonconstruction}
and \ref{sec:ising-production}. Independent state-quality
checks and their error definitions are collected in
Appendix~\ref{sec:method-audit}.

\subsection{What the three supplied papers provide}
Reference~\cite{Herwerth} constructs the $SU(2)_1$ spin chain using vertex
fields and current-mode insertions, including fields at zero and infinity.
Reference~\cite{Montes2016} explains how fusion channels themselves encode
Ising spins, and tests the critical Ising ground-state block.
Reference~\cite{Liu} relates $SU(n)_k$ blocks to projected fermion states:
an auxiliary fermion state is projected onto the allowed on-site spin
representation, facilitating conversion to a matrix product state (MPS).
For $SU(2)_1$ this recovers the projected Fermi-sea/Haldane--Shastry
construction. It does not supply an Ising $\sigma$ eigenstate construction;
$SU(2)_2$ has central charge $3/2$, whereas Ising has $1/2$.

Reference~\cite{Montes2017} relates the homogeneous Ising block vacuum
to the exact fermion pairing state. The covariance reduction is described
in Ref.~\cite{PeschelEisler}. Together these provide the basis for the
Ising production construction in Sec.~\ref{sec:ising-production}. Endpoint-block and projected-state
constructions serve as independent validation and are documented in
Appendix~\ref{sec:ising-independent}.

\subsection{Compact boson: from nonchiral primaries to vertex blocks}
\label{sec:bosonconstruction}
The target theory has independent left and right bosons $\phi_L(z)$ and
$\phi_R(\bar z)$, normalized by logarithmic two-point functions in $z$
and $\bar z$, respectively; mixed left--right contractions vanish.
We use the compactification convention $R=\sqrt2$,
with momentum integer $m$ and winding integer $w$:
\begin{equation}
 Q_L=\frac mR+\frac{wR}{2},\qquad
 Q_R=\frac mR-\frac{wR}{2}.
 \label{eq:bosoncharges}
\end{equation}
Here $R$ fixes the allowed charge lattice. The nonchiral vertex primary is
\begin{equation}
 \mathcal V_{a,b}(z,\bar z)
 ={:}\exp\!\left[\frac{ia}{\sqrt2}\phi_L(z)
                  +\frac{ib}{\sqrt2}\phi_R(\bar z)\right]{:},
 \quad a=m+w,\quad b=m-w.
 \label{eq:fullvertex}
\end{equation}
Colons denote normal ordering, which removes contractions of a field
with itself at the same point. The state created at the origin by this
field is denoted $\ket{as,bs}$, where $s=1/\sqrt2$:
$\ket{as,bs}=\mathcal V_{a,b}(0,0)\ket0$.
This is the state--operator correspondence: a field inserted at the
origin specifies a state on a circle surrounding it.
Each chiral charge $Q$ contributes weight $Q^2/2$, hence
$h=a^2/4$ and $\bar h=b^2/4$. This follows from the stress-tensor
operator product
$T(z)V_Q(0)\sim(Q^2/2)z^{-2}V_Q(0)+z^{-1}\partial V_Q(0)$,
where $T$ generates conformal transformations.

To obtain a finite spin-chain vector we use the auxiliary chiral
endpoint construction of Ref.~\cite{Herwerth}, with the chosen charge
labels at zero and infinity. The resulting chain has a nonchiral
continuum limit; the two endpoint labels specify the target charge
sectors. The auxiliary chiral correlator is the wavefunction amplitude:
we do not multiply it by its complex conjugate to construct a different
state. Its absolute square is taken only when computing probabilities.

The chiral boson is normalized by
$\langle\phi(z)\phi(w)\rangle=-\log(z-w)$. Its vertex field
$V_q(z)={:}\exp(iq\phi(z)){:}$ has weight $q^2/2$. Gaussian contraction gives
\begin{equation}
 \left\langle\prod_r V_{q_r}(z_r)\right\rangle
 =\delta_{\sum_rq_r,0}\prod_{r<t}(z_r-z_t)^{q_rq_t}.
 \label{eq:vertex}
\end{equation}
The neutrality symbol sets the amplitude to zero unless the charges sum to
zero. It fixes the allowed occupation sector, reducing the computation.
More explicitly, a pair of exponentials contributes
$\exp[-q_rq_t\langle\phi(z_r)\phi(z_t)\rangle]
=(z_r-z_t)^{q_rq_t}$, while integration over the spatially constant
boson mode imposes the neutrality condition. Thus the product in
Eq.~\eqref{eq:vertex} is determined by the two-point function.

\begin{enumerate}
\item Put the sites at $z_j=\exp(2\pi i j/N)$, $j=0,\ldots,N-1$.
Associate occupation $n_j$ with a vertex charge $s_j/\sqrt2$, where
$s_j=2n_j-1$. The symbol $s_j$ is a site sign; the state-label symbol
$s=1/\sqrt2$ used below is a continuum charge unit.
\item Insert charges $a/\sqrt2$ at the origin and $b/\sqrt2$ at infinity.
The latter is defined by
$V_q(\infty)=\lim_{W\to\infty}W^{q^2}V_q(W)$.
Neutrality gives $M=\sum_j n_j=(N-a-b)/2$, leaving
$\binom NM$ allowed configurations. These specify the state;
the production contraction in Sec.~\ref{sec:extension} avoids
enumerating them on large chains.
\item Multiply Eq.~\eqref{eq:vertex} by the Marshall sign
$\chi(\bm n)=\exp(i\pi\sum_j jn_j)$, which fixes the spin-basis convention.
After dropping factors independent of $\bm n$, the amplitude is
\begin{equation}
 \psi_{a,b}(\bm n)\propto
 \chi(\bm n)\prod_{i<j}(z_i-z_j)^{s_is_j/2}
 \prod_j(-z_j)^{a s_j/2}
 \prod_j(1-z_j/W)^{b s_j/2}.
 \label{eq:bosonblock}
\end{equation}
Set the last product to one at strict infinity. To reproduce the original
numerics, use $W=10^8$.
\item Evaluate logarithms with one consistent principal branch for each
unordered pair. Subtract the largest real log amplitude before
exponentiating, then normalize. For the occupied set $O$, the implemented
configuration-dependent logarithm is
\begin{align}
 L(O)={}&2\sum_{i<j\in O}\log(z_i-z_j)
 -\sum_{i\in O}\sum_{j\ne i}\log(z_{\min(i,j)}-z_{\max(i,j)})\notag\\
 &+i\pi\sum_{i\in O}i
 +a\sum_{i\in O}\log(-z_i)
 +b\sum_{i\in O}\log(1-z_i/W).
 \label{eq:logblock}
\end{align}
\item Use $(a,b)=(0,0),(1,1),(3,1)$ for the three states. Their
continuum charges and weights are
\begin{center}
\begin{tabular}{lccc}\toprule
State & $(Q_L,Q_R)$ & $(h,\bar h)$ & $M$\\\midrule
$\ket{00}$ & $(0,0)$ & $(0,0)$ & $N/2$\\
$\ket{ss}$ & $(1,1)/\sqrt2$ & $(1/4,1/4)$ & $N/2-1$\\
$\ket{3s,s}$ & $(3,1)/\sqrt2$ & $(9/4,1/4)$ & $N/2-2$\\\bottomrule
\end{tabular}
\end{center}
The different $M$ values ensure exact orthogonality. No orthogonalization
is applied to these states.
\end{enumerate}

The charge-three vertex is a primary of the free-boson Virasoro algebra.
In the extended $SU(2)_1$ current algebra it is a level-two descendant of
the spin-half field, consistent with Ref.~\cite{Herwerth}. These are
compatible uses of ``primary'' for different symmetry algebras.
Algebraically expanding $s_j=2n_j-1$ in Eq.~\eqref{eq:bosonblock} gives the
same Laughlin/Jastrow amplitude as the original generator. We use bit $j$
for site $j$ throughout the production calculation. The explicit numerical
comparison, including conversion of the original basis ordering, is
reported in Appendices~\ref{sec:boson-sanity} and
\ref{sec:state-comparisons}.

\subsubsection{Deriving the Laughlin/Jastrow form}
The following reduction makes the boson construction a directly
reproducible example of the conformal-block method. Let
$O=\{i_1<\cdots<i_M\}$ be the occupied sites and write $d_{ij}=z_i-z_j$
for $i<j$.
\begin{enumerate}
\item Expand the site-charge exponent:
\begin{equation}
 \frac{s_i s_j}{2}=2n_i n_j-n_i-n_j+\frac12.
\end{equation}
The last term contributes the same factor to every configuration and
can be absorbed into normalization. The first term produces
$\prod_{\mu<\nu}(z_{i_\mu}-z_{i_\nu})^2$.
\item Collect the two terms linear in occupation into one factor per
occupied site. Define
\begin{equation}
 D_i=\prod_{j<i}(z_j-z_i)\prod_{j>i}(z_i-z_j).
\end{equation}
Then the configuration-dependent pair product is
$\prod_{\mu<\nu}(z_{i_\mu}-z_{i_\nu})^2\prod_{i\in O}D_i^{-1}$.
The endpoint factors similarly reduce to
$\prod_{i\in O}(-z_i)^a(1-z_i/W)^b$.
\item Use the fact that the $z_j$ are the roots of $z^N-1$:
\begin{equation}
 \prod_{j\ne i}(z_i-z_j)=Nz_i^{N-1},\qquad
 D_i=(-1)^iNz_i^{N-1}.
\end{equation}
The factors $(-1)^i$ cancel the Marshall signs. Since $z_i^N=1$,
$z_i^{-(N-1)}=z_i$. For integer $a$, the remaining
$(-1)^{aM}$ is a configuration-independent phase.
\item The explicit polynomial amplitude is therefore
\begin{align}
 \psi_{a,b}(O)&=\frac{\mathcal P_{a,b}(O;W)}{\sqrt{\mathcal Z_{a,b}(W)}},
 \notag\\
 \mathcal P_{a,b}(O;W)&=
 \prod_{\mu<\nu}(z_{i_\mu}-z_{i_\nu})^2
 \prod_{\mu=1}^{M}z_{i_\mu}^{a+1}(1-z_{i_\mu}/W)^b ,
 \label{eq:bosonpolynomial}
\end{align}
where $\mathcal Z_{a,b}(W)=\sum_{|O|=M}|\mathcal P_{a,b}(O;W)|^2$.
An overall phase is immaterial. Equation~\eqref{eq:bosonpolynomial}
is the exponent-two Laughlin/Jastrow form used by the original method.
At strict infinity, set every $(1-z_i/W)^b$ to one.
\end{enumerate}
Although $b$ disappears from the explicit position factors at infinity,
it still fixes $M=(N-a-b)/2$. Thus changing $b$ changes the state.
The block and polynomial formulas give the same normalized vector,
up to one overall phase; taking an absolute square would destroy this
equivalence.

\subsubsection{Three canonical nonchiral examples}
For each example, the displayed polynomial and its normalization
$\mathcal Z$ specify the amplitudes in the common $2^N$-dimensional spin
basis. The production Pfaffian contraction uses this definition without
storing the full vector. The expressions below are at strict infinity.
\begin{enumerate}
\item \textbf{Vacuum $\ket{00}$.}
Choose $(a,b)=(0,0)$, or momentum/winding $(m,w)=(0,0)$.
The continuum vertex is the identity, with $(h,\bar h)=(0,0)$.
Neutrality requires $M=N/2$, and
\begin{equation}
 \mathcal P_{0,0}(O;\infty)=
 \prod_{\mu<\nu}(z_{i_\mu}-z_{i_\nu})^2\prod_{\mu=1}^{N/2}z_{i_\mu}.
\end{equation}
This is the half-filled Haldane--Shastry ground-state amplitude.
\item \textbf{Vertex primary $\ket{ss}$.}
Choose $(a,b)=(1,1)$, or $(m,w)=(1,0)$. The continuum field is
$\mathopen{:}\exp[i(\phi_L+\phi_R)/\sqrt2]\mathclose{:}$, with
$(h,\bar h)=(1/4,1/4)$, dimension $1/2$, and conformal spin zero.
Neutrality requires $M=N/2-1$, and
\begin{equation}
 \mathcal P_{1,1}(O;\infty)=
 \prod_{\mu<\nu}(z_{i_\mu}-z_{i_\nu})^2
 \prod_{\mu=1}^{N/2-1}z_{i_\mu}^{2}.
\end{equation}
\item \textbf{Vertex primary $\ket{3s,s}$.}
Choose $(a,b)=(3,1)$, or $(m,w)=(2,1)$. The continuum field is
$\mathopen{:}\exp[i(3\phi_L+\phi_R)/\sqrt2]\mathclose{:}$, with
$(h,\bar h)=(9/4,1/4)$, dimension $5/2$, and conformal spin two.
Neutrality requires $M=N/2-2$, and
\begin{equation}
 \mathcal P_{3,1}(O;\infty)=
 \prod_{\mu<\nu}(z_{i_\mu}-z_{i_\nu})^2
 \prod_{\mu=1}^{N/2-2}z_{i_\mu}^{4}.
\end{equation}
Its status as a Virasoro primary follows from the vertex-field
stress-tensor product in Sec.~\ref{sec:bosonconstruction}. Calling it a descendant of the extended
$SU(2)_1$ current algebra refers to a different symmetry algebra.
\end{enumerate}
For the original data convention, multiply the last two polynomials by
$\prod_{\mu}(1-z_{i_\mu}/10^8)$ before normalization; the vacuum is
unchanged. The three vectors have different occupation numbers and
are exactly orthogonal without a Gram--Schmidt procedure. The total
spin component in the convention $S_j^z=s_j/2$ is
$S^z_{\rm tot}=M-N/2=-(a+b)/2$.

The production data use these same normalized vertex/Jastrow states,
with $W=10^8$ fixed to match the original implementation. For larger
chains, their interval matrices are obtained by the exact Pfaffian
contraction in Sec.~\ref{sec:extension}; this preserves the amplitudes
while avoiding full-chain enumeration. The independent polynomial/block
comparison and the Haldane--Shastry sanity check are given in
Appendix~\ref{sec:boson-sanity}.


\subsection{Ising: exact fermion states}
\label{sec:ising-production}
The nonchiral states required by Eqs.~(4.13)--(4.18) of the supplied
manuscript~\cite{theory} are
$I:(0,0)$, $\sigma:(1/16,1/16)$, and $\eps:(1/2,1/2)$, where each
pair gives $(h,\bar h)$. The energy primary $\eps$ contains both
left and right fermions and has scaling dimension one. The production
calculation uses the exact Jordan--Wigner solution for all three states.

The Hamiltonian that identifies these lattice states is
\begin{equation}
 H=-\sum_{j=0}^{N-1}X_jX_{j+1}-\sum_{j=0}^{N-1}Z_j,
 \qquad X_N=X_0,
 \label{eq:Hising}
\end{equation}
with coupling and transverse field both one. These are periodic
\emph{spin} boundary conditions. Define $P=\prod_j Z_j$ and let
$\mathsf T$ translate the spins by one site. The three desired states
have $\mathsf T=1$, corresponding to zero momentum, and are identified by
\begin{center}
\begin{tabular}{lcll}\toprule
State & $P$ & Energy ordering & Production sector\\\midrule
$I$ & $+1$ & Ground state & Antiperiodic vacuum\\
$\eps$ & $+1$ & First excited even state & Two antiperiodic quasiparticles\\
$\sigma$ & $-1$ & Lowest odd state & Periodic vacuum; odd zero mode\\
\bottomrule\end{tabular}
\end{center}
Here $\mathsf T$ is distinct from the stress tensor $T(z)$.

The exact finite-chain energies of these sectors are
\begin{equation}
 E_I=-2\csc\frac\pi{2N},\qquad
 E_\sigma=-2\cot\frac\pi{2N},\qquad
 E_\eps=E_I+8\sin\frac\pi{2N}.
 \label{eq:ising-energies}
\end{equation}
For the velocity $v=2$ of Eq.~\eqref{eq:Hising},
$E_a-E_I\sim(2\pi v/N)(h_a+\bar h_a)$ identifies the spin and energy
primaries with scaling dimensions $1/8$ and $1$. We construct these
sectors directly by their fermion occupations.

\subsubsection{Exact fermion construction used for the production data}
For all three Ising states, the production interval
matrices use the exact Jordan--Wigner solution. It is efficient
because a fermionic Gaussian state is specified by a $2N$ by $2N$
covariance matrix, rather than all $2^N$ amplitudes.
For the spin-to-fermion solution see Ref.~\cite{Montes2017}, Sec.~VIII
and Appendix C; for reduction by correlations see
Ref.~\cite{PeschelEisler}, Sec.~3.3(II), Eqs.~(14)--(18) and the
Majorana discussion immediately following them. The zero-mode parity
selection below fixes the conventions of our implementation explicitly.
\begin{enumerate}
\item \textbf{Transform spins to Majorana fermions.}
Define the Hermitian operators
\begin{equation}
 a_{2j}=\left(\prod_{k<j}Z_k\right)X_j,\qquad
 a_{2j+1}=\left(\prod_{k<j}Z_k\right)Y_j.
\end{equation}
They satisfy $a_ra_s+a_sa_r=2\delta_{rs}$.
In a fixed parity sector,
$H=(i/4)\sum_{r,s}\mathcal A_{rs}a_ra_s$, where the real
antisymmetric matrix has $\mathcal A_{r,r+1}=2$,
$\mathcal A_{r+1,r}=-2$, and boundary entry
$\mathcal A_{2N-1,0}=-2P$.
\item \textbf{Select the correct boundary sector.}
Even spin parity gives antiperiodic fermions, with momenta
$k=(2m+1)\pi/N$. Odd spin parity gives periodic fermions, with
$k=2m\pi/N$. In both cases $m=0,\ldots,N-1$, with momenta understood
modulo $2\pi$. These are two sectors of the same periodic \emph{spin}
Hamiltonian; we have not changed its physical boundary condition.
\item \textbf{Choose the occupation pattern for each primary.}
Diagonalizing the quadratic Hamiltonian defines fermion annihilation
operators $\gamma_k$ with excitation energies
$\omega_k=4|\sin(k/2)|$.
The identity state is their vacuum in the antiperiodic sector:
$\gamma_k\ket{I_N}=0$.
The energy primary occupies the two lowest modes in that sector:
\begin{equation}
 \ket{\eps_N}=\gamma_{\pi/N}^\dagger\gamma_{-\pi/N}^\dagger\ket{I_N}.
\end{equation}
For $\sigma$, use the periodic sector, empty its nonzero-energy modes,
and choose the zero-mode occupation that gives $P=-1$.
The zero mode must be fixed: a covariance with an unresolved zero mode
does not describe the required pure parity eigenstate.
\item \textbf{Construct the covariance in the implemented basis.}
Let $i\mathcal A=U\operatorname{diag}(\lambda_\mu)U^\dagger$.
For the identity, assign mode signs $\nu_\mu=\operatorname{sgn}\lambda_\mu$.
For $\eps$, reverse the signs for the two lowest positive-energy modes
and their negative-energy partners. Then set
\begin{equation}
 G=\operatorname{Re}\!\left[
 iU\operatorname{diag}(\nu_\mu)U^\dagger\right],
 \qquad G_{rs}=\langle ia_ra_s\rangle\quad(r\ne s),\quad G_{rr}=0.
\end{equation}
The mode signs are bookkeeping for the covariance, not occupation
probabilities. In the periodic sector start with $\nu_\mu=0$ on the
zero modes, take real orthonormal null vectors $u,v$ of $\mathcal A$,
and add either $uv^T-vu^T$ or its negative. Choose the sign such that
$(-1)^N\operatorname{Pf}(G)=-1$. The Pfaffian is the antisymmetric-matrix
polynomial obeying $\operatorname{Pf}(G)^2=\det G$; its sign here fixes
the physical fermion parity.
\item \textbf{Check purity and compute the interval matrix.}
Verify $G^T=-G$ and $G^2=-\id$. Compare the energy
$E=\frac14\sum_{r,s}\mathcal A_{rs}G_{rs}$ with
Eq.~\eqref{eq:ising-energies}.
Restrict the covariance to the interval and reconstruct its density
matrix using Eq.~\eqref{eq:gaussian}.
\end{enumerate}
This calculation is performed by \texttt{majorana\_covariance} and
\texttt{gaussian\_rdm}. The numerical purity and full-vector comparisons
are documented in Appendices~\ref{sec:state-comparisons} and
\ref{sec:large-sanity}.


Restrict $G$ to the first $2\ell$ indices. An orthogonal transformation
brings it into $2\times2$ blocks with off-diagonal entries $g_r$. If $b_r$
are the correspondingly rotated Majoranas, the interval state is
\begin{equation}
 \rho_A=\prod_{r=0}^{\ell-1}
 \frac{\id+g_r i b_{2r}b_{2r+1}}2.
 \label{eq:gaussian}
\end{equation}
We build this $2^\ell\times2^\ell$ matrix explicitly. Each component
state is Gaussian; their mixture generally is not. We diagonalize the
actual matrix $m_{ab}$, not an averaged covariance matrix.
For the compact boson, the production Pfaffian contraction evaluates
Eq.~\eqref{eq:rdm} for the fixed-occupation Jastrow states without
explicitly forming the full coefficient matrix. The connected large-chain
procedures for both CFTs are developed in Sec.~\ref{sec:extension}.

% \subsubsection{Why use exact fermion states for production?}
% \label{sec:methodchoice}
% The Ising Hamiltonian is quadratic after the Jordan--Wigner
% transformation. Its covariance construction therefore supplies the
% finite-chain eigenstates directly, with no variational bond-dimension
% truncation or iterative full-vector projection in the production step.
% Fixing the periodic zero-mode parity selects the spin primary, and
% restricting the covariance produces its interval matrix without storing
% $2^N$ amplitudes. This makes the same method applicable across the full
% system-size grid.

% We use the conformal-block and projected vectors as independent checks
% of the state identification, and compare with the supplied periodic
% uniform matrix product states (puMPS). Their construction details are in
% Appendix~\ref{sec:ising-independent}; the numerical overlap, energy,
% residual and reduced-state comparisons are in
% Appendices~\ref{sec:state-comparisons} and
% \ref{sec:ising-spin-validation}. These checks support the production
% choice for this integrable Hamiltonian and the supplied checkpoints;
% they do not establish a general accuracy advantage over fully converged
% puMPS. Exact lattice states still retain finite-size and interval-cutoff
% effects when compared with CFT.
